% Target ~450 words.

Dysarthria is a motor speech disorder caused by neurological injury or
disease --- cerebral palsy, stroke, traumatic brain injury, amyotrophic
lateral sclerosis (ALS), Parkinson's disease --- that affects an estimated
six million people in the United States alone~\cite{Enderby1980Frenchay}.
Severity assessment is the core clinical task: it gates therapy
prescription, tracks progression, and predicts intelligibility outcomes.
Yet assessment today depends on trained Speech-Language Pathologists (SLPs)
applying perceptual instruments such as the Frenchay Dysarthria Assessment
or the Sentence-Intelligibility Test. SLP time is scarce, the rating is
subjective with non-trivial inter-rater variance, and patients in
under-resourced settings often have no access at all.

Automatic severity estimation has therefore been an active research target
for two decades. The dominant paradigm has been to fine-tune a learned
classifier on top of pre-extracted acoustic features --- MFCCs, x-vectors,
or, more recently, the latent codes of self-supervised speech models such
as wav2vec~2.0~\cite{Baevski2020wav2vec2} or
WavLM~\cite{Chen2022WavLM,Hernandez2020SSLDysarthria,Joshi2021IndianDysarthria}.
This paradigm has two structural problems. First, fine-tuning destroys the
pre-trained representation: the very geometry that encodes the regularity
of typical speech is overwritten by the small, noisy supervision signal of
a clinical dataset. Second, it requires labels --- labelled severity ratings
on every patient --- which contradicts the shortage that motivated automation
in the first place.

We propose a complementary approach. WavLM-Large, trained on 94{,}000
hours of neurotypical speech with a denoising objective, has implicitly
learned a manifold of how typical articulators move through phonetic
trajectories. We hypothesise that dysarthric motor impairment manifests
\emph{geometrically} in this manifold: control speakers trace tight,
consistent paths through it, while dysarthric speakers trace erratic,
deviated paths. The geometry \emph{is} the severity signal. Under this
view, severity estimation reduces to measuring the deviation of a
speaker's trajectory from a reference of typical trajectories, and the
SSL model is used \emph{frozen} --- never fine-tuned, never even gradient-
adapted.

We instantiate this idea with Articulatory Trajectory Deviation (ATD), a
length-normalised dynamic-time-warping distance between a speaker's PCA-64
WavLM trajectory and a control-built reference. Three contributions
follow. First, on TORGO, ATD attains speaker-level Spearman
$\rho = 0.756$ ($p = 0.0011$) and a 16.3\% RMSE reduction over MFCC + SVR
baselines, without any fine-tuning. Second, a pre-registered
controlled-perturbation study on neurotypical speech confirms that ATD
responds specifically to articulatory distortion (formant shift, frame
dropout) and not to acoustic factors (time-stretch, additive noise),
ruling out the alternative that ATD merely tracks acoustic distance.
Third, a path-decomposition extension localises deviation to specific
phonemes, providing per-token actionability that scalar metrics lack.
\end{document}% removed; here for refs only
